\section{Project objectives}
\label{sec:project_objectives}
\subsection{Result objectives}
As stated in the project plan (Appendix \ref{app:project_plan}), this project
has a goal of researching, designing and developing a system that as a minimum:

\begin{itemize} 
\item Is more efficient than \bhtb{}.
\item Is developed as an engine that responds to API-calls.
\item Is modular in design and allows future changes and additions that will not
affect existing functionality.
\item Is scalable enough to handle both small and large sets of IDS-sensors and
rules in a production environment.
\item Is able to download rules and rulesets from multiple customizable sources.
\item Is able to distribute different rules to different IDS-sensors.
\item Can download and distribute rules automatically at set intervals.
\item Is capable of exporting rules.
\item Has the capability to activate, deactivate, filter and suppress rules on a
per sensor basis, in a user friendly way.
\item Supports multi-user interaction.
\item Can efficiently and safely work with different installations of \snort{} and various
other third party tools, such as SNORBY \cite{snorby}.
\item Is easy to install.
\item Has an internal logging system.
\item Is well-documented and can be easily adopted by anyone as open source
software.
\end{itemize}
In addition, the project group will attempt to add the following functionality:

\begin{itemize}
\item A graphical user interface.
\item The mentioned user interface must be user friendly and intuitive, and
require only limited knowledge of \snort{} rules from the user.
\item The system can also distribute rules to
Suricata \cite{surikata}, another \snort{}-like IDS, which
can operate with the exact same rules as \snort{}.
\item Implement security features that ensure confidentiality of rules.
\item The project will investigate and potentially implement multi-threading
technology or other optimization techniques to exploit modern hardware
capacities.
\item The system is capable of validating and checking rules for errors before
they are distributed to sensors.
\item The system supports commenting of the rules.
\item The system supports writing custom rules and editing existing rules.
\end{itemize}
Once the system is developed, it will be released as an Open Source Software to
the public in a well-documented format.
\subsection{Effect objectives}

Forming baselines from the current prototype \bhtb{} and from \snortmanager{},
the project has a goal to further increase the efficiency of rule management by
investigating and implementing programming techniques such as multi-threading
and caching to reduce latency and delays in the system by as much as 50\%. The
software must also maintain sublinearity in performance, i.e. 10 executions of
the same job should take 10 times, or less, the amount of time as a single
execution.

The achievement of these goals is to be measured by conducting several
performance tests on the final software solution. Results will also be compared
to the performance of \bhtb{} and \snortmanager{}.
